% Part: turing-machines
% Chapter: undecidability
% Section: enumerating-tms

\documentclass[../../../include/open-logic-section]{subfiles}

\begin{document}

\olfileid{tur}{und}{enu}
\olsection{د ټيورينګ ماشينونو شمېرل}

\begin{explain}
موږ ښودلے شو چې د ټولو ټيورينګ ماشينونو سټ !!{enumerable} دے۔ دا
له دې حقيقت څخه راځي چې هر ټيورينګ ماشين په متناهي ډول بيانېدے شي۔
د حالتونو سټ او د پټې لغت‌زېرمه متناهي سټونه دي۔ د حالت بدلون تابع
له~$Q \times \Sigma$ څخه~$Q \times \Sigma \times \{\TMleft,
\TMright, \TMstay\}$ ته يوه قسمي تابع ده، نو دا هم د هغو متناهي
شمېر آرګوماني جوړو د ارزښتونو په لست کولو ټاکل کېدے شي چې تابع پرې
تعريف شوې ده۔

تر دې ځايه دا خبره سمه ده، خو يو نازک توپير شته۔ د ټيورينګ ماشينونو
تعريف پر دې هېڅ بنديز نۀ لګاوه چې د حالتونو د سټ او د پټې د الفبا
!!{element}s څه شيان کېدے شي۔ نو مثلاً د هر حقيقي عدد دپاره له تخنيکي
پلوه يو ټيورينګ ماشين شته چې هماغه عدد د حالت په توګه کاروي۔ خو د
ټيورينګ ماشين \emph{چلند} له دې څخه خپلواک دے چې کوم څيزونه د حالتونو
او لغت‌زېرمه په توګه کارول کېږي۔ په~\olref{fig:variants} کښې دا دوه
ټيورينګ ماشينونه وګورئ۔
\begin{figure}
\begin{center}
\begin{tikzpicture}[->,>=stealth',shorten >=1pt,auto,node distance=2.8cm,
                    semithick]
  \tikzstyle{every state}=[fill=none,draw=black,text=black]

  \node[initial,state]         (A)              {$q_0$};
  \node[state]         (B) [right of=A] {$q_1$};

  \path (A) edge [bend left] node {\TMtrans{\TMstroke}{\TMstroke}{\TMright}} (B)
        (B) edge [loop above] node {\TMtrans{\TMblank}{\TMblank}{\TMright}} (B)
            edge [bend left] node {\TMtrans{\TMstroke}{\TMstroke}{\TMright}} (A);
\end{tikzpicture}\\
\begin{tikzpicture}[->,>=stealth',shorten >=1pt,auto,node distance=2.8cm,
  semithick]
\tikzstyle{every state}=[fill=none,draw=black,text=black]

\node[initial,state]         (A)              {$s$};
\node[state]         (B) [right of=A] {$h$};

\path (A) edge [bend left] node {\TMtrans{A}{A}{\TMright}} (B)
(B) edge [loop above] node {\TMtrans{\TMblank}{\TMblank}{\TMright}} (B)
edge [bend left] node {\TMtrans{A}{A}{\TMright}} (A);
\end{tikzpicture}
\end{center}
\caption{د \emph{جفت} ماشين بڼې}
\ollabel{fig:variants}
\end{figure}
دا دوه ډياګرامونه له دوو ماشينونو سره سمون لري: يو~$M$ چې د پټې
الفبا ئې~$\Sigma = \{\TMendtape,\TMblank,\TMstroke\}$ او د حالتونو
سټ ئې~$\{q_0,q_1\}$ دے، او بل~$M'$ چې الفبا ئې~$\Sigma' =
\{\TMendtape,\TMblank,A\}$ او حالتونه ئې~$\{s,h\}$ دي۔ خو نورې
لارښوونې ئې يو شان دي: $M$ به د~$n$ دانو~$\TMstroke$ پر لړۍ هله او
يوازې هله ودرېږي چې~$n$ جفت وي، او~$M'$ به د~$n$ دانو~$A$ پر لړۍ
هله او يوازې هله ودرېږي چې~$n$ جفت وي۔ موږ يوازې~$\TMstroke$ په~$A$،
$q_0$ په~$s$، او~$q_1$ په~$h$ بدل نومولي دي۔ دا بېلګه عموميت مومي:
موږ ټيورينګ ماشينونه تر هغه يو شان ګڼلے شو چې يو د نښو او حالتونو د
داسې نوم‌بدلون له لارې له بل څخه تر لاسه کېږي۔ په اصل کښې، موږ د
ټيورينګ ماشين نښې او حالتونه په ساده توګه مثبت صحيح اعداد ګڼلے شو:
د~$\sigma_0$ پر ځاے~$1$، د~$\sigma_1$ پر ځاے~$2$، او داسې نور
وګڼئ؛ $\TMendtape$~$1$ دے، $\TMblank$~$2$ دے، او داسې نور۔ په دې
ډول د \emph{جفت} ماشين په~\olref{fig:standard-even} کښې انځور شوے
ماشين ګرځي۔
\begin{figure}
\[\begin{tikzpicture}[->,>=stealth',shorten >=1pt,auto,node distance=2.8cm,
  semithick]
\tikzstyle{every state}=[fill=none,draw=black,text=black]

\node[initial,state]         (A)              {$1$};
\node[state]         (B) [right of=A] {$2$};

\path (A) edge [bend left] node {\TMtrans{3}{3}{\TMright}} (B)
(B) edge [loop above] node {\TMtrans{2}{2}{\TMright}} (B)
edge [bend left] node {\TMtrans{3}{3}{\TMright}} (A);
\end{tikzpicture}
\]
\caption{يو معياري \emph{جفت} ماشين}
\ollabel{fig:standard-even}
\end{figure}
هغه ټيورينګ ماشين چې حالتونه او نښې ئې مثبت صحيح اعداد وي ښايي
\emph{معياري} ماشين وبولو، او له دې وروسته يوازې معياري ماشينونه
وڅېړو۔\footnote{د ``معياري ماشين'' اصطلاح معياري نۀ ده۔}

موږ غوښتل وښيو چې د ټيورينګ ماشينونو سټ !!{enumerable} دے، او پورتنيو
خبرو ته په پام کافي ده چې وښيو د معياري ټيورينګ ماشينونو سټ
!!{enumerable} دے۔ فرض کړئ يو معياري ټيورينګ ماشين~$M = \tuple{Q,
\Sigma, q_0, \delta}$ راکړل شوے دے۔ هغه د مثبتو صحيح اعدادو په يوه
متناهي مزي څنګه بيانولے شو؟ لومړے به د حالتونو شمېر، پخپله حالتونه،
د نښو شمېر، پخپله نښې، او پيلنی حالت لست کړو۔ (په ياد ولرئ، دا ټول
مثبت صحيح اعداد دي، ځکه~$M$ معياري ماشين دے۔) د~$\delta$ په اړه څه؟
د ممکنو آرګومانونو سټ، يعنې د~$\tuple{q,\sigma}$ جوړې، متناهي دے،
ځکه~$Q$ او~$\Sigma$ متناهي دي۔ نو په~$\delta$ کښې معلومات يوازې د
ټولو هغو $5$-تاييونو~$\tuple{q, \sigma, q', \sigma', d}$ متناهي
لست دے چې~$\delta(q,\sigma) = \tuple{q', \sigma', D}$ وي، او~$d$
هغه عدد دے چې لوری~$D$ کوډوي (مثلاً د~$\TMleft$ دپاره~$1$،
د~$\TMright$ دپاره~$2$، او د~$\TMstay$ دپاره~$3$)۔

په دې ډول هر معياري ټيورينګ ماشين د مثبتو صحيح اعدادو په يوه متناهي
لست، يعنې د~$s_M \in (\PosInt)^*$ په لړۍ، بيانېدے شي۔ مثلاً معياري
\emph{جفت} ماشين په دې لړۍ کوډ کېږي:
\[
2, \underbrace{1, 2}_Q, 3, \overbrace{1, 2, 3}^\Sigma, 1, \underbrace{1, 3, 2, 3, 2}_{\delta(1,3) = \tuple{2,3,R}},
\overbrace{2, 2, 2, 2, 2}^{\delta(2,2) = \tuple{2,2,R}},
\underbrace{2, 3, 1, 3, 2}_{\delta(2,3) = \tuple{1,3,R}}.
\]
\end{explain}

\begin{thm}
له~$\Nat$ څخه~$\Nat$ ته داسې تابعې شته چې د ټيورينګ په وسيله
محاسبه کېدونکې نۀ دي۔
\end{thm}

\begin{proof}
موږ پوهېږو چې د مثبتو صحيح اعدادو د متناهي لړيو سټ~$(\PosInt)^*$
!!{enumerable} دے (\cref{sfr:siz:zigzag:prob:posint-star})۔ له دې څخه
راځي چې د معياري ټيورينګ ماشينونو د بيانونو سټ، د~$(\PosInt)^*$ د
يوه فرعي سټ په توګه، پخپله هم شمېرل کېدونکے دے۔ له~$\Nat$ څخه~$\Nat$
ته هره د ټيورينګ په وسيله محاسبه کېدونکې تابع د کوم (په اصل کښې د
ډېرو) ټيورينګ ماشين له خوا محاسبه کېږي۔ د هغې د حالتونو او نښو نومونه
په مثبتو صحيح اعدادو بدلولو سره (په ځانګړي ډول~$\TMendtape$ په~$1$،
$\TMblank$ په~$2$، او~$\TMstroke$ په~$3$) وينو چې هره د ټيورينګ په
وسيله محاسبه کېدونکې تابع د يوه معياري ټيورينګ ماشين له خوا محاسبه
کېږي۔ معنا دا چې له~$\Nat$ څخه~$\Nat$ ته د ټولو د ټيورينګ په وسيله
محاسبه کېدونکو تابعو سټ هم شمېرل کېدونکے دے۔

له بلې خوا، له~$\Nat$ څخه~$\Nat$ ته د ټولو تابعو سټ !!{enumerable}
نۀ دے (\cref{sfr:siz:red:prob:nat-nat})۔ که ټولې تابعې د کوم ټيورينګ
ماشين په وسيله محاسبه کېدے، نو د هغو ټيورينګ ماشينونو د ټولو بيانونو
په لست کولو مو د ټولو تابعو سټ شمېرلے شو چې دا تابعې محاسبه کوي۔ نو
ځينې تابعې شته چې د ټيورينګ په وسيله محاسبه کېدونکې نۀ دي۔
\end{proof}

\begin{prob}
  ايا د ټيورينګ ماشينونو د بيانولو داسې لار درته ذهن ته راځي چې د
  حالتونو او د الفبا د نښو څرګند لست کولو ته اړتيا ونۀ لري؟ تاسو د
  ``معياري'' ماشين خپل مفهوم تعريفولے شئ، خو د دې په اړه څه ووايئ چې
  ولې هر ټيورينګ ماشين ستاسو په نوې معنا د يوه ``معياري'' ماشين له
  خوا تمثيلېدے شي۔
  \textbf{د ژباړې د نسخې سمون (OLCMP-050):} سرچينې دلته ويل چې هر
  ټيورينګ ماشين د معياري ماشين له خوا ``محاسبه'' کېدے شي؛ ماشين تابع
  نۀ ده، نو دلته موخه د هغه تمثيلول يا مشابه‌چلول دي۔
\end{prob}

\end{document}
